%% das Papierformat zuerst
\documentclass[a4paper, 11pt]{article}
\usepackage[margin=3cm]{geometry}
\usepackage[utf8]{inputenc}
\usepackage[T1]{fontenc}
\usepackage[fleqn]{amsmath} %left aligned equations
\usepackage{hyperref} % clickable refs
\usepackage{graphicx}
\usepackage[toc, numberedsection]{glossaries}
\usepackage{float}
\usepackage{amssymb}
\usepackage{calc}
\usepackage{enumitem} 
\usepackage{url}
\usepackage{parskip}
% http://tex.stackexchange.com/questions/17730/newcommand-and-spacing
\usepackage{xspace}
\usepackage{xparse}
% 'frame' option for figures
\usepackage[export]{adjustbox}
% fancy page headers
\usepackage{fancyhdr}
\usepackage{xcolor}
\usepackage{tabularx}
\usepackage{listings}

% page style
\pagestyle{fancy}
\fancyhead[R]{\includegraphics[width=2cm]{../assets/kitlogo}}
\fancyhead[L]{\leftmark}
\fancypagestyle{plain}{
	\rhead{\includegraphics[width=2cm]{../assets/kitlogo}}
	\lhead{\leftmark}
}

% TODO: template übersetzten

\makeglossaries

%Hack for referencing labels
\makeatletter
\def\namedlabel#1#2{\begingroup
    #2%
    \def\@currentlabel{#2}%
    \phantomsection\label{#1}\endgroup
}
\makeatother
% End: Hack for referencing labels

% Glossar: alle Einträge, aber ohne extra Referenzen
% http://tex.stackexchange.com/questions/115635/glossaries-suppress-pages-when-using-glsaddall
\newcommand*{\glsgobblenumber}[1]{}
\makeatletter
\newcommand*{\glsaddnp}[2][]{
  \glsdoifexists{#2}{
    \def\@glsnumberformat{glsgobblenumber}
    \edef\@gls@counter{\csname glo@#2@counter\endcsname}
    \setkeys{glossadd}{#1}
    \@gls@saveentrycounter
    \@do@wrglossary{#2}
  }
}
\renewcommand{\glsaddallunused}[1][]{
  \edef\@glo@type{\@glo@types}
  \setkeys{glossadd}{#1}
  \forallglsentries[\@glo@type]{\@glo@entry}{
    \ifglsused{\@glo@entry}{}{
      \glsaddnp[#1]{\@glo@entry}}}
}
\makeatother

\renewcommand{\glsnamefont}[1]{\mdseries #1} % glossary entries shouldn’t be bold

% Glossar

% So sieht ein Glossar-Eintrag aus:
%
%\newglossaryentry{dijkstra}{
%  name={Dijkstra’s Algorithmus},
%  description={ein Algorithmus, um den optimalen Pfad in einem gerichteten Graphen zu finden}
%}
%\newglossaryentry{arc}{
%  name={Arc-Flags},
%  description={eine Technik, um Routenberechnung zu beschleunigen},
%  see={dijkstra}
%}
%
% Und so kann er im Dokument verwendet werden:
%
% lorem ipsum dolor sit \gls{arc}, consectetur
%
% End: Glossar

% usage: \counteditem{prefix}{refName} -> item `/prefixXX/` with label `prefix:refName` (where XX is counted in increments of 10)
\makeatletter
\newcommand{\oitem}[2]{
  % define the counter
  \@ifundefined{c@oitem#1}{\newcounter{oitem#1}}{} % initialized at 0
  \addtocounter{oitem#1}{10}
  \item[\namedlabel{#1:#2}{/#1\arabic{oitem#1}/}]
}
\makeatother

% new page after section
\let\oldsection\section
\renewcommand\section{\clearpage\oldsection}

\newcommand{\mamidscreenshot}[1]{\includegraphics[width=\textwidth,frame]{#1}}

\newcommand{\uiel}[3]{\item \textbf{"#1" #2:} #3}

\begin{document}

% place a symbol before clickable links
% this has to come *after* \begin{document} because hyperref installs a \AtBeginDocument hook that updates the ref command.
\newcommand{\refsymbol}[0]{\scalebox{0.5}{$\nearrow$}}
\let\oldref\ref
\renewcommand{\ref}[1]{\refsymbol\oldref{#1}}
\let\oldgls\gls
\renewcommand{\gls}[1]{\refsymbol\oldgls{#1}}
\let\oldGls\Gls
\renewcommand{\Gls}[1]{\refsymbol\oldGls{#1}}
\let\oldglspl\glspl
\renewcommand{\glspl}[1]{\refsymbol\oldglspl{#1}}
\let\oldGlspl\Glspl
\renewcommand{\Glspl}[1]{\refsymbol\oldGlspl{#1}}
\let\oldglslink\glslink
\renewcommand{\glslink}[2]{\refsymbol\oldglslink{#1}{#2}}
\let\oldhyperref\hyperref
\renewcommand{\hyperref}[2][notActuallyOptional]{\refsymbol\oldhyperref[#1]{#2}}
\let\oldautoref\autoref
\renewcommand{\autoref}[1]{\refsymbol\oldautoref{#1}}

\newcommand{\abbildung}[1]{\autoref{fig:#1}}
\newcommand{\mamid}{\textit{MAMID}\xspace}

\NewDocumentCommand{\gotype}{m m g g}{
  \subsection{#1}
  #2
  \IfValueT{#4}{
    \subsubsection{Properties}
    \begin{description}
      #4
    \end{description}
   }
   \IfValueT{#3}{
      \subsubsection{Methods}
      \begin{description}
        #3
      \end{description}
    }
}
\NewDocumentCommand{\struct}{m m g g}{\gotype{Struct #1}{#2}{#4}{#3}}
\newcommand{\interface}[3]{\gotype{Interface #1}{#2}{#3}}

\NewDocumentCommand{\property}{m o m+}{
  \item{\textbf{#1\IfValueT{#2}{: #2}} #3}
}
\NewDocumentCommand{\method}{m o d() m+}{
  \item[]{\textbf{#1\IfValueT{#2}{: #2}}}
  \IfValueT{#3}{
    \begin{description}
      #3 % should only contain \param
    \end{description}
  }
  #4  
}
\NewDocumentCommand{\param}{m d<> m}{
  \item{\textbf{\IfValueT{#2}{#2: }#1} #3}
}

% alle Glossareintraege
\newacronym{gui}{GUI}{Graphical User Interface}
\newacronym{cli}{CLI}{Command Line Interface}
\newacronym{CRUD}{CRUD}{Create / Read / (Update|Modify) / Delete}
\newacronym{ICMP}{ICMP}{Internet Control Message Protocol}
\newacronym{API}{API}{Application Programming Interface}
\newacronym{JSON}{JSON}{JavaScript Object Notation}
\newacronym{HTTP}{HTTP}{HyperText Transfer Protocol}
\newacronym{LAN}{LAN}{Local Area Network}
\newacronym{WAN}{WAN}{Wide Area Network}

\newglossaryentry{cluster}{
	name={cluster},
	description={Aggregation of \glspl{host}},
	plural={clusters}
}
\newglossaryentry{MongoDB}{
	name={MongoDB},
	description={A NoSQL based database server},
	plural={MongoDB}
}
\newglossaryentry{replica set}{
	name={replica set},
	description={Multiple \gls{MongoDB} instances distributed over multiple \glspl{host} sharing a copy of the same data},
	plural={replica sets}
}
\newglossaryentry{sharding}{
	name={Sharding},
	description={A type of \gls{MongoDB} deployment where data sets are spread across multiple \glspl{host} or \glspl{replica set}. \\ More information: {\small \url{https://docs.mongodb.com/manual/sharding/}}}
}
\newglossaryentry{administrator}{
	name={administrator},
	description={The person responsible for managing the hardware \& software deployed on the \gls{cluster}. Uses \mamid for \gls{MongoDB} \gls{replica set} administration},
	plural={administrators}
}
\newglossaryentry{host}{
	name={host},
	description={Network node with a unique hostname},
	plural={hosts}
}
\newglossaryentry{slave}{
	name={slave},
	description={Program running on \gls{host} managing \gls{MongoDB} processes},
	plural={slaves}
}
\newglossaryentry{master}{
	name={master},
	description={Program responsible for managing \glspl{slave} and providing \acrshort{API} for \gls{cluster} management},
	plural={masters}
}
\newglossaryentry{inventory}{
	name={inventory},
	description={Persistently stored list of slaves and their state (see \ref{D:Inventory})},
	plural={inventories}
}
\newglossaryentry{active mode}{
	name={active mode},
	description={Marks the \gls{slave} as available to host \gls{MongoDB} processes that are be part of \glspl{replica set}.},
	plural={active modes}
}
\newglossaryentry{maintenance mode}{
	name={maintenance mode},
	description={Marks the \gls{slave} as unavailable \& inhibits reconfiguration through \gls{master}, but does not otherwise affect running \gls{MongoDB} instances on the \gls{host}.},
	plural={maintenance modes}
}
\newglossaryentry{disabled mode}{
	name={disabled mode},
	description={A \gls{slave} in this mode does not run any \gls{MongoDB} instances controlled by the \gls{master}, hence has no \gls{MongoDB} instances in any \gls{replica set}.},
	plural={disabled modes}
}
\newglossaryentry{risk group}{
	name={risk group},
	description={A set of \glspl{host} sharing a common risk of failure, e.g. a shared power supply. Modeled through sets of \glspl{slave} since a 1:1 relationship exists between hosts and slaves},
	plural={risk groups}
}
\newglossaryentry{persistent storage}{
	name={persistent storage},
	description={Storage capable of storing data between power outage},
	plural={persistent storages}
}
\newglossaryentry{volatile storage}{
	name={volatile storage},
	description={Storage incapable of storing data between process lifetime},
	plural={volatile storages}
}
\newglossaryentry{root data directory}{
	name={root data directory},
	description={Base directory holding all files of a \gls{slave} and its \gls{MongoDB} instances},
	plural={root data directories}
}
\newglossaryentry{arbiter}{
	name={arbiter},
	description={In charge of breaking ties on a \gls{replica set} election with even number of members. More information: {\small \url{https://docs.mongodb.com/manual/core/replica-set-elections/}}},
	plural={arbiters}
}
\newglossaryentry{degraded}{
	name={degraded},
	description={State of a \gls{replica set} with member in \gls{maintenance mode} or \gls{disabled mode} or otherwise reporting failure},
	plural={degraded}
}

\begin{titlepage}
\makeatletter
\begin{center}
~\\[4em]
{\Huge MAMID}\\[.8em]\huge{Monitor and Manager for In-memory Databases}\\[2em]
{\huge Design}\\[1em]
{\large\today}\\[2.5em]
{\LARGE
Niklas Fuhrberg\\
Anton Schirg\\
Christian Schwarz\\
Janis Streib\\
Bob Weinand\\[3em]}
{\Large supervised by}\\[2em]
{\LARGE
Dr Marek Szuba\\[1em]}
{\Large at}\\[1em]
{\LARGE
Karlsruhe Institute of Technology\\
SCC\\[2em]}
{\color{gray}
  \small Document Version: \newtheorem{theorem}{Theorem}


\section{master}

\renewcommand{\gocurpackage}{model}
\renewcommand{\gocurpackage}{master}

\struct{ClusterAllocator}{
  The \refstruct{ClusterAllocator} determines the layout of the cluster managed by \mamid.

  It attempts to fulfill the constraints defined through the model objects, in particular
  \begin{itemize}
    \item A \refgo{model.ReplicaSet}'s \refgo{VolatileNodeCount} \& \refgo{PersistentNodeCount}
    \item The \refgo{model.Slave}'s allowed number of Mongod instances
          (\refgo{MongodPortRangeBegin} to \refgo{MongodPortRangeEnd}).
    \item The \refgo{model.Slave.SlaveState}
    \item The configured \refgo{model.RiskGroup}s.
  \end{itemize}

  An iterative algorithm is employed to decide on a cluster layout described through
  \refgo{model.Mongod.DesiredState}s that
  \begin{itemize}
    \item attempts to fulfill the above constraints
    \item attempts an even distribution of Mongods on the different cluster hosts
    \item is a minimal change in comparison to the previous layout
  \end{itemize}

  \begin{theorem}{Idempotence of the ClusterAllocator}
    \label{theorem:idempotence_clusterallocator}
    Let $l$ be a layout of the cluster. Then $ClusterAllocator(ClusterAllocator(l)) = ClusterAllocator(l)$.
  \end{theorem}

}{
  \property{DB}[gorm.DB]{Initialized handle to the database.}
  \property{BusChannel}[chan interface{}]{Initialized channel to the application bus.} %TODO ref}
}{
  \method{LayoutCluster}{Lay out the cluster as described above.}
}

\subsubsection{Pseudocode}

The \refstruct{master.ClusterAllocator} is crucial to the stable operation of the \mamid-managed cluster.\\
Hence, it is worth defining the implementation of \refgo{master.ClusterAllocator.LayoutCluster} through pseudocode. %TODO parentheses

While studying the algorithms below, the reader should keep in mind that
\begin{itemize}
  \item changes in the cluster layout $\equiv$ change or creation of \refgo{model.Mongod.DesiredState} 
  \item the \refgo{model.Mongod.ObservedState} may change after an arbitrary amount or even never.\\
  \item changes to a ReplicaSet must not violate or further worsen the high-availibility constraints,
        in particular \refgo{model.ReplicaSet}'s \refgo{VolatileNodeCount} \& \refgo{PersistentNodeCount}\\
        $\implies$ Mongods in \refgo{model.MongodExecutionState.Recovering} are an important special-case.
\end{itemize}

\newcommand{\pluseq}{\mathrel{{+}{=}}}
\newcommand{\minuseq}{\mathrel{{-}{=}}}
\SetKwInOut{Input}{input}
\SetKwInOut{Output}{output}
\SetKwProg{Fn}{Function}{}{}
\SetKw{Invariant}{invariant}
\IncMargin{0.5em}

Invariant: Algorithm(Algorithm(x)) = Algorithm(x)

Otherwise oscillations could occur

Most interesting case: Set a slave to disabled. Then a new Mongod should be spawned and recover and when it is done the old disabled one can be deleted.

\begin{algorithm}

\caption{Count members of a Replica Set that are \& intended to be in stable state (running)}

\Input{ReplicaSet $r$}
\Output{Number of $p_e$ and $v_e$ member processes of $r$ that are fully operational and are planned to remain in that state.}
\BlankLine
\Fn{EffectiveMemberCount(r ReplicaSet)}{

$p_e, v_e = 0$

\For{$m \in \text{r.Mongods}$}{
	\If{$\text{m.ObservedState.ExecutionState} = Running$ \\ 	%	#this line evaluates to false if m.ObservedState = NULL
		$\land \text{m.DesiredState.ExecutionState} = Running$}{ %#this line evaluates to false if m.DesiredState = NULL 

		\uIf{$m.ParentSlave.PersistentStorage$}{
			$p_e \pluseq 1 $
		}\Else{
			$v_e \pluseq 1 $
		}
	}
}
\Return $p_e, v_e$
}
\end{algorithm}


\begin{algorithm}
\caption{Destroy members of disabled replica sets where possible without violating p/v constraints.}
\ForEach{r in ReplicaSets}{
	
	$p_e, v_e \gets $ EffectiveRunningMembers(r)
	
	\ForEach(// same for persistent and volatile){$\textbf{x}_e \in \{p_e, v_e\}$}{
		\While{$x_e > r.\textbf{x}$}{
			\textbf{destroy} any $m \in r.Mongods$ where $m.ParentSlave$ is $\textbf{x} \land \textbf{disabled}$\;
			$\textbf{x}_e \minuseq 1$
		}
		
		\Invariant minimum number of disabled slaves are running members of $r$\;
		
		\While{$\textbf{x}_e > r.\textbf{x}$}{
			\textbf{destroy} any $m \in r.Mongods$ where $m.ParentSlave$ is $\textbf{x}$\;
			$\textbf{x}_e \minuseq 1$
		}
		
		\Invariant at most $r.\textbf{x}$ members 
	}
	
	\Invariant \textit{desired state}: at most (r.p|r.v) member processes of $r$\;
	// desired state = the state that will be deployed
	
}
\end{algorithm}

\begin{algorithm}
\caption{Count recovering and active members of a Replica Set}

\Input{ReplicaSet $r$}
\Output{Number of $p_a$ and $v_a$ member processes of $r$ that \begin{itemize}
		\item are \textbf{recovering}, i.e. soon-to-be fully operational
		\item fully operational (\textbf{running})
	\end{itemize} and should actually be in that state.}
\BlankLine
\Fn{AlreadyAddedMemberCount(r ReplicaSet)}{
	$p_a, v_a = 0$
	
	\For{$m \in \text{r.Mongods}$}{
		\If{$\text{m.ParentSlave.ConfiguredState} != Disabled$ \\
			$\land \text{m.DesiredState.ExecutionState} != NotRunning$}{ 
			
			\uIf{$m.ParentSlave.PersistentStorage$}{
				$p_a \pluseq 1 $
			}\Else{
			$v_a \pluseq 1 $
		}
	}
}
\Return $p_a, v_a$
}
\end{algorithm}

\begin{algorithm}
\caption{Spawn Mongods on under-provisioned Replica Sets respecting RiskGroup \& p/v constraints.}

$p_a, v_a \gets $ AlreadyAddedMemberCount(r)

\ForEach(// same for persistent and volatile){$\textbf{x}_a \in \{p_a, v_a\}$}{

	%# TODO need to fill the queue. only ReplicaSets which 'need' (use p_* to define what need means) are in  the queue 
	$RQ \gets$ PriorityQueue(R)\;% "relative amount of missing members")\;
	$RGSQ \gets$ map[RiskGroup]PriorityQueue(Slaves of RiskGroup)\;%, "relative amount of available Mongod ports")\;
	
	\While{$r = RQ.pop()$; $r \neq nil$}{
	
		\Invariant $r$ needs a member of type $\textbf{x}$
		
		// find $x$-slave $s$ in a risk group $g$ with no other Mongod of $r$ in $g$\;
		$candidates \gets \{g \mid g \in RGSQ.keys \land \text{r has no member in g}\}$\;
		$g \gets$ ARGMAX($g \in candidates$, RGSQ[g].PeekPriority())\;
		\uIf{$g \neq nil$}{
			$s \gets RGSQ[g].pop()$\;
			\Invariant $s$ has at least one free port\;
			spawn new Mongod $m$ on $s$ and add it to $r.Mongods$\;
			compute $MongodState$ for $m$ and set the $DesiredState$ variable\;
			
			\If{$r$ needs another member of type $\textbf{x}$}{
				$RQ.insert(r)$ // recompute priority
			}
			
			\If{$s$ has free ports}{
				$RGSQ[g].insert(s)$ // recompute priority
			}
			
			
		} \Else{
			yield error\;
			continue\;
		}
		
	
	}

}

\end{algorithm}


}
\end{center}
\makeatother
\end{titlepage}
\newpage
\tableofcontents
\newpage


%http://tex.stackexchange.com/questions/43002/how-to-preserve-the-same-parskip-in-minipage
\newlength{\currentparskip}
\newenvironment{minipageparskip}
{\setlength{\currentparskip}{\parskip}% save the value
	\begin{minipage}{\textwidth}% open the minipage
		\setlength{\parskip}{0.8em}% restore the value
	}
	{\end{minipage}}

% Args: Method, URI, 
\newcommand{\apicall}[8][]{
\begin{minipageparskip}
	\paragraph{\framebox{\uppercase{#1}} #2 - #3}
	\begin{description}[leftmargin=!,labelwidth=\widthof{\bfseries Return Codesaa}] %aa intentioal -  more distance
		\item[Parameters]
		\if\relax\detokenize{#4}\relax
		\emph{None}
		\else
		\begin{tabularx}{\linewidth}{c|c|*1{>{\centering\arraybackslash}X}@{}}
			\textbf{Parameter} & \textbf{Data Type} & \textbf{Description}\\
			\hline
			#4
		\end{tabularx}
		\fi
		
		\item[Returns]
		\begin{description}[leftmargin=!,labelwidth=\widthof{\bfseries Status Codeaa}] %aa intentioal -  more distance
			\item[Status Code] #5
			\ifthenelse{\equal{#6}{}}
			{} % optional argument #1 is empty: skip
			{\item[Headers] \begin{tabularx}{\linewidth}{c|*1{>{\centering\arraybackslash}X}@{}}
					\textbf{HTTP Header} & \textbf{Content}\\
					\hline
					#6
			\end{tabularx}}
			\ifthenelse{\equal{#7}{}}
			{\item[Body] \emph{Empty}} % optional argument #1 is empty: skip
			{\item[Body] #7}
		\end{description}
		
		\item[Error Codes]
		\begin{tabularx}{\linewidth}{c|*1{>{\centering\arraybackslash}X}@{}}
			\textbf{HTTP Status Code} & \textbf{Reason}\\
			\hline%
\ifthenelse{\equal{#1}{put}}{400 & Invalid parameters in object}{}%
\ifthenelse{\equal{#1}{get}}{404 & Object does not exist}{}%
\ifthenelse{\equal{#1}{post}}{404 & Object does not exist\\400 & Invalid parameters in object or parameter may not be changed}{}%
\ifthenelse{\equal{#1}{delete}}{404 & Object does not exist}{}%
\if\relax\detokenize{#8}\relax
\else
\\#8
\fi
		\end{tabularx}
	\end{description}
\end{minipageparskip}
\vspace{1em}
}

\section{Master API}
%TODO General api description
%Uses RESTful HTTP
%Generic error responses
\subsection{API Objects}
API Objects are exchanged in JSON format.
\subsubsection{Slave}
\begin{lstlisting}
{
	"id": <int>,
	"hostname": <string>,
	"slave_port": <int>,
	"mongod_port_range_begin": <int (inclusive)>,
	"mongod_port_range_end": <int (exclusive)>,
	"persistant_storage": <bool>,
	"root data directory": <string>,
	"state": <string>
}
\end{lstlisting}
\subsubsection{Replica Set}
\begin{lstlisting}
{
	"id": <int>,
	"name": <string>,
	"persistent_node_count": <int>,
	"volatile_node_count": <int>,
	"configure_as_sharding_config_server": <bool>
}
\end{lstlisting}
\subsubsection{Risk Group}
\begin{lstlisting}
{
	"id": <int>,
	"name": <string>
}
\end{lstlisting}
\subsection{API Methods}
\subsubsection{Slaves}
%Do not add default error codes
\apicall[get\relax]{/slaves}{Gets all slaves}
	{}
	{200}
	{}
	{array of slave objects}
	{}

\apicall[put]{/slave}{Adds a slave}
	{slave & slave object & slave to add}
	{201}
	{Location & Reference to added slave}
	{}
	{}
	
\apicall[get]{/slave/<id>}{Gets a slave}
	{id & integer & id of the slave to get}
	{200}
	{}
	{slave object}
	{}
	
\apicall[post]{/slave/<id>}{Updates a slave}
	{id & integer & id of the slave to update}
	{200}
	{}
	{}
	{}
	
\apicall[delete]{/slave/<id>}{Deletes a slave}
	{id & integer & id of the slave to delete}
	{200}
	{}
	{}
	{}
	
\subsubsection{Replica Sets}
%Do not add default error codes
\apicall[get\relax]{/replicasets}{Gets all replica sets}
	{}
	{200}
	{}
	{array of replica set objects}
	{}

\apicall[put]{/replicaset}{Adds a replica set}
	{replicaset & replica set object & replica set to add}
	{201}
	{Location & Reference to added replica set}
	{}
	{}

\apicall[get]{/replicaset/<id>}{Gets a replica set}
	{id & integer & id of the replica set to get}
	{200}
	{}
	{replica set object}
	{}

\apicall[post]{/replicaset/<id>}{Updates a replica set}
	{id & integer & id of the replica set to update}
	{200}
	{}
	{}
	{}

\apicall[delete]{/replicaset/<id>}{Deletes a replica set}
	{id & integer & id of the replica set to delete}
	{200}
	{}
	{}
	{}

\apicall[get]{/replicaset/<id>/slaves}{Gets a replica sets slaves}
	{id & integer & id of the replica set to get the slaves of}
	{200}
	{}
	{array of slave objects}
	{}

\subsubsection{Risk Groups}
%Do not add default error codes
\apicall[get\relax]{/riskgroups}{Gets all risk groups}
	{}
	{200}
	{}
	{array of risk groups objects}
	{}

\apicall[put]{/riskgroup}{Adds a risk group}
	{risk group & risk group object & risk group to add}
	{201}
	{Location & Reference to added risk group}
	{}
	{}

\apicall[get]{/riskgroup/<id>}{Gets a risk group}
	{id & integer & id of the risk group to get}
	{200}
	{}
	{risk group object}
	{}

\apicall[post]{/riskgroup/<id>}{Updates a risk group}
	{id & integer & id of the risk group to update}
	{200}
	{}
	{}
	{}

\apicall[delete]{/riskgroup/<id>}{Deletes a risk group}
	{id & integer & id of the risk group to delete}
	{200}
	{}
	{}
	{}
	
\apicall[put]{/riskgroup/<risk\_group\_id>/slave/<slave\_id>}{Assigns a slave to a risk group}
	{risk\_group\_id & integer & id of the risk group to assign the slave to\\
	 slave\_id & integer & id of the slave to assign to the risk group}
	{200}
	{}
	{}
	{}

\apicall[delete]{/riskgroup/<risk\_group\_id>/slave/<slave\_id>}{Removes a slave from a risk group}
	{risk\_group\_id & integer & id of the risk group to remove the the slave from\\
	slave\_id & integer & id of the slave to remove from to the risk group}
	{200}
	{}
	{}
	{}
\section{Slave}
The slave is installed on the individual host nodes and started by the main service manager of the operating system during the boot process.

\subsection{Application / main()}{
  A Go application starts by initializing the Go 'main' package and then running the \textit{main()} function.
  This function performs early initialization of the slave's main datastructures:
  \begin{itemize}
    \item Parse command line flags using an external CLIFlagsParser.
    \item Initialize the MSPListener.
    \item Initialize the Controller with the MSPListener instance.
    \item Transfer control to MSPListener to wait for incoming connections.
  \end{itemize}
}

\struct{MSPListener}{
  Structure for handling communication through the MasterSlaveProtocol.
}{
  \property{consumer}{MSPConsumer}{
    Reference to an event handler, typically a controller.
  }
}{
  \method{Run}{
    Start listening for connections using the MasterSlaveProtocol.
    \begin{itemize}
      \item Listen for incoming connections from the Master.
      \item Decode / Encode the transport format.
      \item Validate incoming request format.
      \item Delegate valid requests to the MSPConsumer.
    \end{itemize}
  }
}

\interface{MSPConsumer}{
  An MSPConsumer handles MasterSlaveProtocol requests received by an instance of \ref{slave:MSPListener}.
}{
  \method{MSPStatusRequest}[SlaveStatus]{
    %TODO reference request message object here -> proto specification
    Request to provide a status report of the slave's state.
    See \ref{msp:SlaveStatus} for details on the data returned by this method.
  }
  \method{MSPEstablishMongodState}[error](
    \param{m}{MSP.Mongod}{State description of a mongod instance that shall be established}
  ){
    %TODO reference request message object here -> proto specification
    Request to establish the configuration state $m$ of a MongoDB process on the slave.
    May return an error in case the state could not be established.
  }
}

\struct{Controller}{
  Handles MSP requests by implementing the MSPConsumer interface. %TODO ref
  Hence, it coordinates the work required to fulfill requests from the master.
  Most importantly, it leverages ProcessManager and MongodConfigurator to spawn and configure instances of MongoDB.
}{
    \property{BusyTable}{map[PortNumber]bool}{
    Contains a boolean busy state per instance of MongoDB to prevent concurrent configuration of a single instance.
  }
  \property{busyTableLock}{sync.Mutex}{
    Mutex controlling access to busyTable. %TODO ref
  }
}

\struct{ProcessManager}{
  Starts processes, may provide a list of alive processes and eventually kill these.
}{
  \property{runningProcesses}{map[PortNumber]*exec.Cmd}{
    Holds the active process controls per instance
  }
}{
  \method{spawnProcess}[error](
    \param{m}[MSP.MongoD]{Mongod information about what exactly to spawn}
    \param{dataDir}[string]{Root directory of MongoDB data}
  ){
    Spawns a Mongod process as requested by the given Mongod inside the data root directory
  }
  \method{runningProcesses}[[]PortNumber]{
    Returns the PortNumbers of the currently running processes
  }
  \method{killProcess}[error](
    \param{p}[PortNumber]{PortNumber to identify process}
  ){
    Kills process by given PortNumber. Does not error if the process is already dead, only if it could not be killed.
  }
  \method{killProcesses}[error]{
    Kills all remaining Processes. Errors if it could not kill some processes.
  }
}

\interface{MongodConfigurator}{
  Applies or returns configuration of a mongod instance by PortNumber.
}{
  \method{MongodConfiguration}[{(MSP.MongodConfiguration, error)}](
    \param{p}{PortNumber}{PortNumber to connect to locally}
  ){
    Reads configuration from a local MongoDB instance
  }
  \method{ApplyMongodConfiguration}[error](
    \param{p}{PortNumber}{PortNumber to connect to locally}
    \param{c}{MSP.MongodConfiguration}{Configuration to apply}
  ){
    Connects to a local MongoDB instance on the given port and applies the configuration
  }
}

\struct{ConcreteMongodConfigurator}{
  Implements MongodConfigurator interface and uses an external MongoDB client to communicate with local instances
}


\glsaddallunused
\makeatletter
\newglossarystyle{myAltlist}{
  \glossarystyle{altlist} % base this style on altlist
  \renewcommand*{\glossaryentryfield}[5]{
  \item[\glsentryitem{##1}\glstarget{##1}{##2}]
    \mbox{}\par\nobreak\@afterheading
    ##3\glspostdescription\space On page ##5.
  }
}
\makeatother
\printglossary[type=main, title={Glossary}, toctitle={Glossary}, style=myAltlist]

\end{document}
